% Generated draft: requires mathematical and source review.
\documentclass[11pt]{article}
\usepackage[margin=25mm]{geometry}
\usepackage{fontspec}
\setmainfont{Latin Modern Roman}
\usepackage{amsmath,amssymb,mathtools,unicode-math}
\setmathfont{Latin Modern Math}
\usepackage{enumitem,needspace,xurl,hyperref}
\setlength{\parindent}{0pt}
\setlength{\parskip}{5pt}
\setlength{\emergencystretch}{3em}
\newcommand{\N}{\mathbb{N}}
\newcommand{\Z}{\mathbb{Z}}
\newcommand{\Q}{\mathbb{Q}}
\newcommand{\R}{\mathbb{R}}
\newcommand{\C}{\mathbb{C}}
\newcommand{\F}{\mathbb{F}}
\newcommand{\A}{\mathbb{A}}
\newcommand{\PP}{\mathbb{P}}
\newcommand{\E}{\mathbb{E}}
\newcommand{\eps}{\varepsilon}
\newcommand{\dd}{\,\mathrm{d}}
\begin{document}
\section*{UnsolvedMath ID: 30005923}
\textbf{Source problem code:} catalogue entry
\subsection{Definitions and mathematical statement}
Using the standard complex-matrix and probability-Haar conventions, let \(M_k(\C)\) denote the complex \(k\times k\) matrices, and define the row norm of \(A=(A_1,\ldots,A_d)\in M_k(\C)^d\) by
\[
\|A\|_{\mathrm{row}}:=\left\|\sum_{j=1}^d A_jA_j^*\right\|^{1/2},
\]
where \(^*\) denotes conjugate transpose and the norm is the operator norm on \(\C^k\). For each positive integer \(N\), let
\[
U(N):=\{V\in M_N(\C):V^*V=VV^*=I_N\},
\]
and equip \(U(N)^d\) with product normalized Haar measure \(\dd\mu_N^{\otimes d}(U)\). For \(U=(U_1,\ldots,U_d)\), put
\[
A\otimes U:=\sum_{j=1}^d A_j\otimes U_j\in M_{kN}(\C),
\]
where the tensor products on the right are Kronecker products. These standard conventions make explicit notation left undefined in the supplied excerpt; their agreement with the report requires review.

\textbf{Conjecture 12 (supplied formulation).} For every real \(r\) with \(0\le r<1\) and every integer \(k\ge1\), there exists a finite constant \(C(r,k)\ge0\) such that, for every integer \(d\ge1\) and every \(A\in M_k(\C)^d\) satisfying \(\|A\|_{\mathrm{row}}\le r\),
\[
\sup_{N\in\{1,2,\ldots\}}
\int_{U(N)^d}
\left|\det\!\left(I_{kN}-\sum_{j=1}^d A_j\otimes U_j\right)\right|^2
\,\dd\mu_N^{\otimes d}(U)
\le C(r,k).
\]
In particular, the same constant must work independently of \(d\), \(N\), and \(A\).

\textbf{Source qualifications.} The excerpt presents this bound as a weaker conjecture whose proof would make rigorous a formal argument, inspired by a strong Szeg\H{o} limit theorem proof attributed to Bump and Diaconis, for a preceding conjecture. The included Conjecture 10 concerns polynomials \(p,q\) in noncommuting indeterminates \(x_1,\ldots,x_d\), with \(p(0)=q(0)=1\), assumed stable and admitting determinantal representations \(\det p(X)=\det(I-A\otimes X)\) and \(\det q(X)=\det(I-B\otimes X)\). Its transcribed limit formula has unclear conjugation, and stability is not defined in the supplied record, so that contextual formula is not reconstructed here.

The excerpt reports that a related conjecture holds for all \(d\) when \(k=1\), and for all \(k\) when \(d=1\); the referent of that sentence is unclear because both Conjectures 12 and 10 appear in the transcription. It also proposes a particular bound involving \(1-r^2\), but the exponent is corrupted by a detached numeral \(2\), so no precise candidate constant is asserted here. The source further suggests that, for fixed \(d\) and \(A\), these integrals should be nondecreasing in \(N\), and reports that both the proposed explicit bound and this monotonicity hold when \(k=1\). These are supplied contextual claims, not additional requirements of the displayed conjecture or independently verified facts.
\subsection{Short English statement}
For each \(0\le r<1\) and matrix size \(k\ge1\), the conjecture asks for a finite bound on the squared determinant averaged over independent Haar unitary matrices, uniform in their size \(N\), the number \(d\) of matrices, and every tuple with row norm at most \(r\). Standard notation and measure conventions are made explicit above; source ambiguities and reported special cases remain subject to human review.
\subsection{Sources}
\begin{itemize}
\item Local UnsolvedMath registry, permanent ID \(30005923\), title: \emph{Uniform Determinant Integrals for Row Contractions}; exact source problem code: \texttt{catalogue entry}. The supplied record labels the problem open and the statement \texttt{corrected\_verified}, with a supplied claim that it was reconstructed and checked against the official report. These labels and that verification claim have not been independently checked here. No individual proposer is supplied; the proposed year is 2024.
\item \emph{Non-commutative Function Theory and Free Probability} (2024), Oberwolfach Reports, pp.\ 1227--1284, DOI 10.4171/owr/2024/22: \url{https://doi.org/10.4171/owr/2024/22}. The supplied excerpt identifies the target as Conjecture 12, equation (13), and includes contextual material from Conjecture 10, equation (11).
\item Supplied literature source: \emph{Noncommutative Szego-type theorem}, \url{https://arxiv.org/abs/2507.19453}. No authors or further bibliographic details are supplied in the record.
\item Literature assessment dated exactly \texttt{2026-08-22}: the supplied curation labels the problem open, reports that the uniform determinant integral bound is known in special cases, and states that no general proof or counterexample for arbitrary row contractions and dimensions was found. This is supplied dated curation, not an independently verified literature assessment; its status and cited claims require independent verification.
\item The excerpt attributes a proof of the strong Szeg\H{o} limit theorem to Bump and Diaconis, under reference label [1]. The corresponding bibliographic entry is absent from the supplied record.
\end{itemize}
\end{document}
